\chapter{Build: an Agent-Based Market}\label{ch:mx:build-an-agent-based-market}

An \omterm{def:mx:build-a-matching-engine-and-exchange-simulator:sim}{exchange simulator} with nobody on it is an empty room. Put a few kinds of simple traders in it, some who know nothing, some who know the value, some
who follow trends and one who makes markets, and it starts to look like a market; which of them it needs in order to look like one is itself a finding.
This chapter describes the agent library this book has used since chapter~11, \texttt{firm.agentmkt}, adds the two agent types it lacked (chartists and
an inventory-skewing market maker), calibrates the population to a target market by the \omterm{def:mx:build-an-agent-based-market:msm}{method of simulated moments}, and then takes it apart one agent
type at a time to see which fact each one carries.

\section{Agents and their information}

\begin{definition}[Agent-based model]\label{def:mx:build-an-agent-based-market:abm}
An \emph{agent-based model}\index{agent-based model} of a market is a simulation in which prices and volumes are not specified but emerge from the
orders of many simulated traders (agents), each following its own rule with its own information, through a trading mechanism such as a \omterm{def:mx:the-limit-order-book:lob}{limit order
book}.
\end{definition}

\begin{definition}[Fundamentalist agent, chartist agent]\label{def:mx:build-an-agent-based-market:types}
A \emph{fundamentalist agent}\index{fundamentalist agent} trades toward a value it believes the asset is worth: it buys when the price is below that
value and sells when it is above, more often the larger the gap. A \emph{chartist agent}\index{chartist agent} trades on the price's own recent
history, typically in the direction of its recent trend.
\end{definition}

Agent-based markets come in two traditions. The first gives agents strategies and asks what dynamics follow: Lux and Marchesi (1999) built a market of
fundamentalists and chartists whose news process had neither fat tails nor any time dependence in volatility, and found that the interactions of the
agents produced both. The second removes strategy almost entirely: Farmer, Patelli and Zovko (2005) let agents place and cancel orders at random
(chapter~1's \omterm{def:mx:the-limit-order-book:zi}{zero-intelligence model}) and, on eleven London Stock Exchange stocks, explained 96\% of the cross-sectional variance of the spread and 76\% of
that of the price diffusion rate with one free parameter. Chiarella and Iori (2002) put heterogeneous rule-based traders into an order-driven market with
a central matching mechanism and studied how their strategies, the tick and the lifetime of orders shape spreads, volume and volatility. The population
here takes a little from each: most of its order flow is random, and a few rules give it direction.

\begin{omfigure}
\begin{tikzpicture}[font=\footnotesize,box/.style={draw,rounded corners=2pt,minimum height=0.6cm,inner sep=3pt,align=center,text width=3.5cm}]
\node[box,fill=gray!12,text width=2.6cm,minimum height=1.4cm] (x) at (0,0) {\textbf{exchange simulator}\\ \texttt{firm.exchsim}\\ one book, one feed};
\node[box,fill=omDef!12] (lp) at (-4.9,1.5) {liquidity providers\\ limit orders, cancels};
\node[box,fill=omDef!12] (nz) at (-4.9,0) {noise takers\\ market orders, sign runs};
\node[box,fill=omDef!12] (fd) at (-4.9,-1.5) {fundamentalists\\ see the hidden value $V$};
\node[box,fill=omProp!14] (ch) at (4.9,1.5) {chartists\\ follow the recent trend};
\node[box,fill=omProp!14] (mm) at (4.9,0) {market maker\\ skews on inventory};
\node[box,fill=omThm!14] (ex) at (4.9,-1.5) {execution agents\\ metaorders, algorithms};
\foreach \n in {lp,nz,fd} {\draw[->,thick] (\n.east) -- (x.west);}
\foreach \n in {ch,mm,ex} {\draw[->,thick] (\n.west) -- (x.east);}
\node[anchor=north] at (0,-2.3) {one Population agent: independent Poisson clocks, zero latency};
\draw[dashed,omDef] (-6.8,-2.25) rectangle (-3.0,2.2);
\end{tikzpicture}

\omcaption{The population. The three types on the left run as one agent, \texttt{Population}, on independent Poisson clocks; the chartists belong to it
too, but are new in this chapter; the market maker and the execution agents are separate agents with their own sessions and a 20-microsecond
latency.}\label{fig:mx:build-an-agent-based-market:pop}
\end{omfigure}

\texttt{Population} runs its traders as clocks whose rates are recomputed after every event
(\cref{fig:mx:build-an-agent-based-market:pop}). Liquidity providers add one-lot \omterm{def:mx:the-limit-order-book:orders}{limit orders} on each side, half of them one tick from the opposite
best quote and the rest up to ten ticks behind it, and every resting order is cancelled at a constant rate. Noise takers send \omterm{def:mx:the-limit-order-book:orders}{market orders} whose signs
come in runs (next section). Fundamentalists send \omterm{def:mx:the-limit-order-book:orders}{market orders} toward a hidden value $V$, a random walk of one-tick jumps, at a rate proportional to
the distance between $V$ and the mid: they are also this market's \omterm{def:mx:why-there-is-a-spread:informed}{informed traders} (chapter~4), since nobody else sees $V$. Chartists send \omterm{def:mx:the-limit-order-book:orders}{market
orders} in the direction of the mid's change over the last 30 seconds, at a rate proportional to its size. The noise orders and the jumps of $V$ are
drawn in advance from the seed, so two runs that differ by one agent share them (common random numbers).

\omcode{code/firm/agentmkt/firm_agentmkt.py}{159}{187}{The population's event: an exogenous noise order or jump of the hidden value, or else one of its clocks chosen in proportion to its rate, a provider's limit order, a cancellation, a fundamentalist's or a chartist's market order.}\label{lst:mx:build-an-agent-based-market:pop}

\section{Liquidity providers and their inventory}

\begin{definition}[Market-making agent]\label{def:mx:build-an-agent-based-market:mm}
A \emph{market-making agent}\index{market-making agent} quotes on both sides continuously and manages the inventory its fills leave it with, usually by
moving its quotes against that inventory so as to attract the trades that reduce it (inventory skew).
\end{definition}

The population's providers never hold anything: their orders are independent of what they have filled. The new \texttt{MarketMaker} does: every half
second it quotes two lots at the other traders' best bid and ask (one tick inside when their spread is wider than two ticks), both moved against its
inventory by a skew in ticks per lot held, and it stops adding to a position beyond 20 lots. Its reference excludes its own quotes; an early version
that did not walked its own bid down, refresh after refresh, whenever it was long, since each new quote was set below the previous one.

\omcode{code/firm/agentmkt/firm_agentmkt.py}{251}{272}{The market maker's refresh: the others' best quotes, a shift of the skew times the inventory in lots, a quote kept if it is already at the right price and replaced otherwise.}\label{lst:mx:build-an-agent-based-market:mm}

\begin{table}[ht]
\centering
\small
\begin{tabular}{@{}lrrr@{}}
\toprule
skew (ticks per lot) & inventory (lots, root mean square) & profit per hour (USD) & spread (ticks) \\
\midrule
0 & 14.5 & $-116$ (203) & 1.031 \\
0.1 & 4.2 & 448 (98) & 1.023 \\
0.3 & 1.6 & 545 (67) & 1.024 \\
\bottomrule
\end{tabular}
\caption{The market maker in the calibrated market at three inventory skews, three hours each: the time-weighted root mean square of its inventory,
its profit (cash after fees plus the inventory at the last mid) with the standard deviation across hours, and the market's mean spread. Data:
\texttt{mx\_agents.maker\_study}.}\label{tab:mx:build-an-agent-based-market:mm}
\end{table}

Without skew the maker's inventory wanders (\cref{fig:mx:build-an-agent-based-market:inv}) and it loses money: its fills follow the fundamentalists,
who know the value, so the inventory it accumulates is the wrong way round when $V$ moves (chapter~4's adverse selection). With a skew of 0.3 ticks per
lot its inventory stays within a few lots and its hourly profit is positive and steadier (\cref{tab:mx:build-an-agent-based-market:mm}); the skew is
the practical form of the inventory models of chapter~4 (and of Book~2, chapter~30). The market's spread hardly changes: the providers keep it near one tick
whatever the maker does.

\begin{omfigure}
\begin{tikzpicture}
\begin{axis}[width=10.6cm,height=5.4cm,xlabel={\small seconds from the open},ylabel={\small inventory (lots)},
  tick label style={font=\footnotesize},grid=major,grid style={gray!25},xmin=0,xmax=3600,
  legend columns=2,legend style={font=\footnotesize,at={(0.5,-0.3)},anchor=north,draw=none,fill=none,
  /tikz/every even column/.append style={column sep=8pt}},table/col sep=comma]
\addplot[omThm,thick,const plot] table[x=t,y=lots]{figdata/microstructure/27-build-an-agent-based-market/inventory_skew0.csv};
\addplot[omDef,thick,const plot] table[x=t,y=lots]{figdata/microstructure/27-build-an-agent-based-market/inventory_skew3.csv};
\legend{no skew,skew 0.3 ticks per lot}
\end{axis}
\end{tikzpicture}

\omcaption{The market maker's inventory through one hour of the calibrated market, with the same order flow, without skew and with a skew of 0.3 ticks
per lot. Without skew it swings from one 20-lot limit to the other. Data: \texttt{mx\_agents.maker\_study}.}\label{fig:mx:build-an-agent-based-market:inv}
\end{omfigure}

\section{Execution agents and metaorders}

A metaorder (Book~7, chapter~9) cut into children sends many \omterm{def:mx:the-limit-order-book:orders}{market orders} of one sign; this is the mechanism Lillo, Mike and Farmer (2005) proposed for
the long memory of order signs: if the sizes of large orders have a power-law tail, $P(V > v) \propto v^{-\alpha}$, and they are executed in children of
constant size, the autocorrelation of order signs decays as $\tau^{-(\alpha - 1)}$, which is long memory (not summable) when $\alpha < 2$. The noise
takers of \texttt{firm.agentmkt} are that mechanism in a line: their signs come in runs whose lengths are Pareto with tail $\alpha$ (the parameter
\texttt{run\_tail}), each run a metaorder. Execution agents proper, \texttt{Metaorder} for a schedule of market-order children and the algorithms of
chapters 16 to 22, are added to the same simulator as separate agents; chapter~11 measured their impact against the same session without them.

\begin{omfigure}
\begin{tikzpicture}
\begin{loglogaxis}[width=9.6cm,height=5.6cm,xlabel={\small lag (orders)},ylabel={\small sign autocorrelation},
  tick label style={font=\footnotesize},grid=major,grid style={gray!25},xmin=0.8,xmax=120,ymin=0.001,ymax=0.5,
  legend columns=2,legend style={font=\footnotesize,at={(0.5,-0.3)},anchor=north,draw=none,fill=none,
  /tikz/every even column/.append style={column sep=6pt}},table/col sep=comma,unbounded coords=discard]
\addplot[omDef,very thick,mark=*,mark size=1.3pt] table[x=lag,y=runs]{figdata/microstructure/27-build-an-agent-based-market/signs.csv};
\addplot[gray,very thick,mark=square*,mark size=1.3pt] table[x=lag,y=none]{figdata/microstructure/27-build-an-agent-based-market/signs.csv};
\addplot[omProp,very thick,mark=triangle*,mark size=1.6pt] table[x=lag,y=excess]{figdata/microstructure/27-build-an-agent-based-market/signs.csv};
\addplot[omThm,thick,dashed] table[x=lag,y=law]{figdata/microstructure/27-build-an-agent-based-market/signs.csv};
\legend{with sign runs,without,excess from the runs,slope $-(\alpha-1)=-1.5$}
\end{loglogaxis}
\end{tikzpicture}

\omcaption{Autocorrelation of the signs of aggressive orders in the calibrated market (run tail $\alpha = 2.5$), with and without the noise takers'
sign runs, three hours each; the excess the runs add, with the power law their tail implies. The excess is drawn out to lag 20, where it is lost in noise. Data:
\texttt{mx\_agents.sign\_curves}.}\label{fig:mx:build-an-agent-based-market:signs}
\end{omfigure}

The calibrated market has $\alpha = 2.5$, so its runs predict an excess autocorrelation falling as $\tau^{-1.5}$: short memory. Measured over lags 1 to
7 the excess falls with a log-log slope of $-1.15$ (\cref{fig:mx:build-an-agent-based-market:signs}), a little flatter than the asymptotic law, and is
gone by lag 10. What remains without the runs, about 0.07 to 0.09 out to lag 20, is the fundamentalists' doing: while $V$ is away from the mid they
send order after order in the same direction.

\section{Calibrating to stylised facts}

\begin{definition}[Method of simulated moments]\label{def:mx:build-an-agent-based-market:msm}
The \emph{method of simulated moments}\index{method of simulated moments} estimates the parameters of a model that can be simulated but whose moments
have no closed form: it chooses the parameters that minimise a weighted distance between moments computed from the data and the same moments averaged
over simulations of the model, the simulations using the same random draws at every candidate parameter.
\end{definition}

Franke and Westerhoff (2012) estimated small fundamentalist-chartist models this way, with moments chosen from the stylised facts (Book~7, chapter~5)
of daily index returns. Here the target is one hour of Book~7's synthetic market, \texttt{firm.tape} without its news window, measured on eight hours:
it has Hawkes-clustered noise flow, metaorders, \omterm{def:mx:why-there-is-a-spread:informed}{informed traders} and a random activity level, none of which the agent population has as such. The
seven facts are the excess kurtosis of 30-second mid returns, their lag-1 autocorrelation, volatility clustering (the mean autocorrelation of absolute
returns at lags 1 to 5), the autocorrelation of the signs of aggressive orders at lags 1 and 10 (the prints of one order counted once), the
time-weighted mean spread and the mean depth at the best. The distance is the sum of squared differences between simulated and target means, each
scaled by the target's standard deviation across its eight hours.

\omcode{code/firm/agentmkt/firm_agentmkt.py}{307}{322}{The method of simulated moments on a grid: the distance in target standard deviations, and the average of the same seeds' facts at every grid point.}\label{lst:mx:build-an-agent-based-market:msm}

Three parameters are calibrated, on a grid of three values each: the noise takers' rate (0.9, 1.2, 1.5 a second), their run tail (none, 1.5, 2.5) and
the chartists' rate (0, 0.05, 0.1 a second per tick of trend); the providers, fundamentalists and market maker (skew 0.1) are fixed. Each point is run
for three hours with the same three seeds.

\begin{omfigure}
\begin{tikzpicture}
\begin{axis}[width=9.2cm,height=5.4cm,xlabel={\small chartists' rate (per second per tick of trend)},ylabel={\small distance},
  tick label style={font=\footnotesize},grid=major,grid style={gray!25},xtick={0,0.05,0.1},
  xticklabel style={/pgf/number format/fixed},ymin=0,
  legend columns=3,legend style={font=\footnotesize,at={(0.5,-0.3)},anchor=north,draw=none,fill=none,
  /tikz/every even column/.append style={column sep=6pt}},table/col sep=comma]
\addplot[gray,very thick,mark=square*,mark size=1.4pt] table[x=chart,y=tail0]{figdata/microstructure/27-build-an-agent-based-market/surface.csv};
\addplot[omThm,very thick,mark=triangle*,mark size=1.6pt] table[x=chart,y=tail15]{figdata/microstructure/27-build-an-agent-based-market/surface.csv};
\addplot[omDef,very thick,mark=*,mark size=1.4pt] table[x=chart,y=tail25]{figdata/microstructure/27-build-an-agent-based-market/surface.csv};
\legend{no sign runs,tail 1.5,tail 2.5}
\end{axis}
\end{tikzpicture}

\omcaption{The distance between simulated and target facts at a noise rate of 1.2 orders a second, by the chartists' rate and the run tail, three hours
per point with the same seeds. Data: \texttt{mx\_agents.calibration}.}\label{fig:mx:build-an-agent-based-market:surface}
\end{omfigure}

The best point is a noise rate of 1.2, a run tail of 2.5 and a chartists' rate of 0.05, at a distance of 10.1 over its three calibration hours; the
starting point (noise 0.9, no runs, no chartists) was at 27.9 (\cref{fig:mx:build-an-agent-based-market:surface}). On six fresh hours the calibrated
population is at 5.4, and six fresh hours of the target market itself are at 1.7 from the target's mean: what sampling noise alone gives.
\Cref{fig:mx:build-an-agent-based-market:facts} shows where the rest comes from: the calibrated market has too little volatility clustering and too
little depth, and a little too much sign memory at lag 1.

\begin{omfigure}
\begin{tikzpicture}
\begin{axis}[width=10.6cm,height=5.6cm,ybar,bar width=7pt,ylabel={\small simulated $-$ target (target s.d.)},
  xtick={0,1,2,3,4,5,6},xticklabels={kurtosis,return acf,clustering,sign acf 1,sign acf 10,spread,depth},
  x tick label style={font=\footnotesize,rotate=25,anchor=north east},tick label style={font=\footnotesize},
  grid=major,grid style={gray!25},enlarge x limits=0.08,
  legend columns=2,legend style={font=\footnotesize,at={(0.5,-0.42)},anchor=north,draw=none,fill=none,
  /tikz/every even column/.append style={column sep=8pt}},table/col sep=comma]
\addplot[fill=gray!45,draw=gray] table[x=k,y=start]{figdata/microstructure/27-build-an-agent-based-market/facts.csv};
\addplot[fill=omDef!55,draw=omDef] table[x=k,y=calibrated]{figdata/microstructure/27-build-an-agent-based-market/facts.csv};
\legend{starting point (calibration hours),calibrated (six fresh hours)}
\end{axis}
\end{tikzpicture}

\omcaption{Each fact's gap to the target in target standard deviations, for the starting point and for the calibrated population. The squares of the
calibrated bars sum to the distance on fresh hours. Data: \texttt{mx\_agents.calibration}.}\label{fig:mx:build-an-agent-based-market:facts}
\end{omfigure}

The ablation takes the calibrated population and removes one agent type at a time on the six fresh hours, with the same seeds: the sign runs (the noise
takers keep their rate but their signs become independent), the chartists, the fundamentalists, the market maker. \Cref{tab:mx:build-an-agent-based-market:abl}
gives each fact's move in target standard deviations.

\begin{table}[ht]
\centering
\small
\begin{tabular}{@{}lrrrrrrr@{}}
\toprule
removed & kurtosis & return & clustering & sign & sign & spread & depth \\
 & & acf & & acf 1 & acf 10 & & \\
\midrule
sign runs & $+1.32$ & $+0.02$ & $-0.63$ & $\mathbf{-2.77}$ & $-0.35$ & $-0.38$ & $+0.10$ \\
 & (0.79) & (0.29) & (1.12) & (0.15) & (0.25) & (0.19) & (0.06) \\[2pt]
chartists & $\mathbf{+1.11}$ & $-0.10$ & $-0.72$ & $+0.02$ & $-0.18$ & $-0.29$ & $+0.09$ \\
 & (0.55) & (0.20) & (0.67) & (0.10) & (0.27) & (0.11) & (0.02) \\[2pt]
fundamentalists & $\mathbf{+8.55}$ & $-0.47$ & $-0.73$ & $+0.87$ & $-2.42$ & $-1.29$ & $+0.96$ \\
 & (2.65) & (0.25) & (0.88) & (0.14) & (0.40) & (0.15) & (0.05) \\[2pt]
market maker & $-0.79$ & $-1.77$ & $+0.31$ & $-0.69$ & $-1.00$ & $\mathbf{+9.44}$ & $-0.75$ \\
 & (0.29) & (0.25) & (0.57) & (0.10) & (0.19) & (0.37) & (0.04) \\[2pt]
\bottomrule
\end{tabular}
\caption{Ablation: the change in each fact when one agent type is removed from the calibrated population, in target standard deviations, the mean of
six hours with the same seeds (paired standard errors in brackets below; the largest move of each type in bold). Data:
\texttt{mx\_agents.ablation}.}\label{tab:mx:build-an-agent-based-market:abl}
\end{table}

Each type carries something (\cref{tab:mx:build-an-agent-based-market:abl}). The sign runs carry the lag-1 sign memory ($-2.77$, standard error 0.15 standard deviations).
The fundamentalists carry the price's movement: without them the mid is pinned between rare jumps, so the kurtosis of returns explodes ($+8.6$) and the
long-lag sign memory falls ($-2.42$). The market maker carries the spread ($+9.4$): without it, the spread widens by 0.15 ticks. The chartists
carry nothing this test can see: their largest move, on kurtosis, is $+1.11$ with a standard error of 0.55, two standard errors. Removing the noise takers altogether leaves a market in
which only fundamentalists trade, with a sign autocorrelation near one; it is not in the table because every fact moves.

\section{What agent-based markets can and cannot tell you}

An agent-based market is a tool for questions about mechanisms: what happens to spreads if the tick halves (chapter~6), whether a band attracts the
price (chapter~25), how an algorithm's shortfall depends on the other traders' reaction (chapter~28). Replaying a historical tape cannot answer them,
because a replayed market does not react; Vyetrenko and coauthors (2020) compared market replay with interactive agent-based configurations on a list
of realism metrics taken from the stylised facts, and found the interactive ones more realistic, and more so when the agents' fundamental value came from
historical data rather than a mean-reverting random walk. ABIDES (Byrd, Hybinette and Balch 2020), an open-source simulator of tens of thousands of
agents with an exchange agent whose messages follow Nasdaq's ITCH and OUCH, is built for exactly this.

What it cannot tell you is why real traders do what they do. Matching seven facts to within 5.4 does not identify the agents: other populations,
with other rules, would match them as well (here the chartists could be removed at little cost), and the facts left unmatched, volatility clustering
first, say that something in the target (its random activity level) is missing from the population. The calibration answers ``is this population good
enough for the question?'', not ``is this how the market works?''; and the answer depends on the facts chosen and their weights. The fair use is
comparative: the same population, with and without an algorithm, a rule or a trader, and the same seeds.

\section{Tutorial: a market that looks like a market}\label{tut:mx:build-an-agent-based-market:tut}

\begin{tutorial}
\textbf{Goal.} Populate \texttt{firm.exchsim} with agents, calibrate them to a target market's stylised facts, and find which fact each agent type
carries. \textbf{End state:} \cref{tab:mx:build-an-agent-based-market:mm,tab:mx:build-an-agent-based-market:abl},
\cref{fig:mx:build-an-agent-based-market:inv,fig:mx:build-an-agent-based-market:signs,fig:mx:build-an-agent-based-market:surface,fig:mx:build-an-agent-based-market:facts}.
\begin{tutsteps}
\item \textbf{Populate.} \texttt{PopulationConfig(lo\_rate, near, cancel, depth, noise, run\_tail, fund, v\_rate, chart, chart\_window)};
\texttt{session(cfg, seconds, seed, agents=[MarketMaker(skew=0.1)])}; the \texttt{Result}'s \texttt{tape()} is in \texttt{firm.tape}'s format.
\item \textbf{Measure.} \texttt{facts(tape, 30.0)}; chapter~3's \texttt{firm\_lobstats.report(tape)} for the book's own facts (depth profile,
lifetimes, \omterm{def:mx:empirical-facts-of-order-books:life}{fleeting orders}, resilience).
\item \textbf{Calibrate.} \texttt{mx\_agents.target()}, \texttt{calibration()} (\texttt{msm} over the grid with common seeds).
\item \textbf{Take apart.} \texttt{ablation()}, \texttt{sign\_curves()}, \texttt{maker\_study()}; draw with \texttt{fig\_agents.py}.
\end{tutsteps}
\textbf{What to change next.} Add a switching rule between fundamentalists and chartists (Lux and Marchesi's mechanism) and see whether the clustering
gap closes; target a real tape's facts; weight the moments by their simulated covariance instead of the target's variances.
\end{tutorial}

\section{Build: agentmkt}\label{bld:mx:build-an-agent-based-market:agentmkt}

\begin{build}
\bfield{Purpose} The agent library and population runner on \texttt{firm.exchsim}, used since chapter~11 and by the tutorials of Books 11 and 12: a
reacting market with known ingredients.
\bfield{Interface} \texttt{PopulationConfig}, \texttt{Population(cfg, seed, seconds)} (\texttt{.exo}, \texttt{.v\_path}), \texttt{Metaorder(plan)},
\texttt{MarketMaker(lots, skew, max\_lots, refresh\_s)}, \texttt{session(cfg, seconds, seed, agents, venue)}; \texttt{facts(tape, sample\_s)},
\texttt{FACTS}, \texttt{distance(sim, target, scale)}, \texttt{msm(run, grid, target, scale, seeds)}.
\bfield{Rules} Deterministic and seeded per agent; exogenous flow drawn in advance so that runs share it; output in \texttt{firm.tape}'s format; the
default population unchanged by later additions (the chartists and the activity curve are off by default).
\bfield{Acceptance tests} \texttt{code/firm/agentmkt/tests/}: the same seed gives the same tape and the exogenous flow is shared; a metaorder trades its
plan; the activity curve shapes the flow; the market maker only rests and bounds its inventory; one sign per aggressive order; \texttt{msm} finds the
minimum of a known function.
\bfield{Stretch} Strategy switching; agents with learning rules; several instruments and venues in one population.
\end{build}

\begin{omsources}
\item T.~Lux and M.~Marchesi, ``Scaling and criticality in a stochastic multi-agent model of a financial market'', \emph{Nature} 397, 1999.
\item J.~D.~Farmer, P.~Patelli and I.~I.~Zovko, ``The predictive power of zero intelligence in financial markets'', \emph{Proceedings of the National
Academy of Sciences} 102(6), 2005.
\item C.~Chiarella and G.~Iori, ``A simulation analysis of the microstructure of double auction markets'', \emph{Quantitative Finance} 2(5), 2002.
\item F.~Lillo, S.~Mike and J.~D.~Farmer, ``Theory for long memory in supply and demand'', \emph{Physical Review E} 71, 066122, 2005.
\item R.~Franke and F.~Westerhoff, ``Structural stochastic volatility in asset pricing dynamics: estimation and model contest'', \emph{Journal of
Economic Dynamics and Control} 36(8), 2012.
\item D.~Byrd, M.~Hybinette and T.~H.~Balch, ``ABIDES: towards high-fidelity multi-agent market simulation'', Proceedings of the 2020 ACM SIGSIM
Conference on Principles of Advanced Discrete Simulation.
\item S.~Vyetrenko, D.~Byrd, N.~Petosa, M.~Mahfouz, D.~Dervovic, M.~Veloso and T.~H.~Balch, ``Get real: realism metrics for robust limit order book
market simulations'', Proceedings of the First ACM International Conference on AI in Finance, 2020.
\end{omsources}

\section{Exercises}

\begin{exercise}[$\star$]\label{exo:mx:build-an-agent-based-market:1}
The target's lag-1 sign autocorrelation is 0.158 with a standard deviation of 0.046 across its hours; the calibrated population gives 0.212 on fresh
hours. How many target standard deviations is the gap, and what does it add to the distance?
\end{exercise}

\begin{exercise}[$\star$]\label{exo:mx:build-an-agent-based-market:2}
By Lillo, Mike and Farmer's result, at what rate does the sign autocorrelation decay for run tails of 1.5 and 2.5? Which is long memory?
\end{exercise}

\begin{exercise}[$\star$]\label{exo:mx:build-an-agent-based-market:3}
If the population were exactly right and the facts' simulated variances equal to the target's, what distance would you expect from the mean of six
simulated hours against the mean of eight target hours, with seven facts? Compare with the fresh distance and the floor.
\end{exercise}

\begin{exercise}[$\star\star$]\label{exo:mx:build-an-agent-based-market:4}
Why is the calibrated point's distance on its three calibration hours larger than on six fresh hours?
\end{exercise}

\begin{exercise}[$\star\star$]\label{exo:mx:build-an-agent-based-market:5}
Why does removing the fundamentalists make the kurtosis of returns explode?
\end{exercise}

\begin{exercise}[$\star\star$]\label{exo:mx:build-an-agent-based-market:6}
Why does the market maker lose money without inventory skew and make money with it, when its quotes are otherwise the same?
\end{exercise}

\begin{exercise}[$\star\star\star$]\label{exo:mx:build-an-agent-based-market:7}
\emph{Coding.} Recalibrate on the five return and order-flow facts only (drop spread and depth), from the grid's stored facts. Does the best point
change, and what is its distance?
\end{exercise}

\begin{exercise}[$\star\star\star$]\label{exo:mx:build-an-agent-based-market:8}
\emph{Find the flaw.} ``Our calibrated agent-based market reproduces the stylised facts, so real traders are fundamentalists, chartists and \omterm{def:mx:why-there-is-a-spread:informed}{noise
traders} in these proportions.''
\end{exercise}

\section{Problem: A Market That Looks Like a Market}

\begin{problem}[{Weekend problem --- a market that looks like a market}]\label{pb:mx:build-an-agent-based-market:1}
A research lead wants a simulated market to test \omterm{def:mx:benchmark-algorithms:algo}{execution algorithms} in, and asks how realistic it is and what makes it so.

\medskip\noindent\textbf{Part I --- Agents.}
\begin{enumerate}
\item Define an \omterm{def:mx:build-an-agent-based-market:abm}{agent-based model}, a \omterm{def:mx:build-an-agent-based-market:types}{fundamentalist agent} and a \omterm{def:mx:build-an-agent-based-market:types}{chartist agent}.
\item What did Lux and Marchesi find, and what did Farmer, Patelli and Zovko find?
\item Describe the population's liquidity providers, noise takers, fundamentalists and chartists.
\item Why are the fundamentalists also the \omterm{def:mx:why-there-is-a-spread:informed}{informed traders} here?
\item Why are the noise orders drawn in advance from the seed?
\end{enumerate}

\medskip\noindent\textbf{Part II --- Inventory and metaorders.}
\begin{enumerate}[resume]
\item Define a \omterm{def:mx:build-an-agent-based-market:mm}{market-making agent} and describe \texttt{MarketMaker}'s quotes.
\item What went wrong when the maker's reference included its own quotes?
\item Give the maker's inventory and profit at the three skews.
\item State Lillo, Mike and Farmer's result, and how the noise takers implement it.
\item What does the sign-autocorrelation figure show about the calibrated tail?
\end{enumerate}

\medskip\noindent\textbf{Part III --- Calibration.}
\begin{enumerate}[resume]
\item Define the \omterm{def:mx:build-an-agent-based-market:msm}{method of simulated moments}, and say why the same seeds are used at every grid point.
\item What is the target, and what are the seven facts?
\item Give the grid and the best point.
\item \emph{State the named result}: the distance between simulated and target facts after calibration, in and out of sample, against the starting
point and the floor, and the fact each agent type carries.
\item Which facts does the calibrated population miss, and why?
\end{enumerate}

\medskip\noindent\textbf{Part IV --- Judgement.}
\begin{enumerate}[resume]
\item What did Vyetrenko and coauthors find about replay and agent-based markets?
\item Why can a replayed tape not answer an execution question that an agent market can?
\item Why does matching the facts not identify the agents?
\item What would you tell the research lead about using this market to compare algorithms?
\item In one sentence: what is an agent-based market for?
\end{enumerate}
\end{problem}

\section{Interview questions}

\begin{interviewq}[$\star$ \iqroles{researcher}]\label{iq:mx:build-an-agent-based-market:1}
What is a zero-intelligence market, and what can it explain?
\end{interviewq}

\begin{interviewq}[$\star\star$ \iqroles{researcher}]\label{iq:mx:build-an-agent-based-market:2}
How would you calibrate an agent-based market to a stock, and how would you know you had overfitted?
\end{interviewq}

\begin{interviewq}[$\star\star$ \iqroles{developer}]\label{iq:mx:build-an-agent-based-market:3}
How do you make an agent-based simulation deterministic, and why does it matter for comparing two algorithms?
\end{interviewq}

\begin{interviewq}[$\star\star$ \iqroles{trader}]\label{iq:mx:build-an-agent-based-market:4}
Why should a market maker skew its quotes on inventory, and by how much?
\end{interviewq}

\begin{interviewq}[$\star\star$ \iqroles{mle}]\label{iq:mx:build-an-agent-based-market:5}
You train an execution policy by reinforcement learning in an agent-based market. What can go wrong when it meets the real one?
\end{interviewq}

\begin{interviewq}[$\star\star\star$ \iqroles{risk}]\label{iq:mx:build-an-agent-based-market:6}
A desk validates a new algorithm only in an agent-based simulator. What would you ask before approving it?
\end{interviewq}
